\section{Conclusions}
\label{sec:conclusion}

Public practitioner discussions can support digital-laboratory research when they are converted into bounded, auditable objects instead of counted as field incidents. The Observatory represents that conversion through explicit source records, coding rules, field-level metric inputs, robustness analyses, claim controls, schemas, executable code, and versioned releases.

The empirical analysis identifies ten recurrent bottlenecks and a common structure. Operational difficulty concentrates where a digital representation must support a physical or scientific commitment. Interfaces require accessible execution paths. Source requires deployment identity. Resource values require physical validation. Event streams require command--state--material linkage. Recovery requires final disposition. AI-generated methods require evidence at the claimed validation stage.

The strongest prospective mechanism links observability and recovery. Real incident datasets can now test whether linked events predict safe-state time, salvage, disposition completeness, and recurrence. The released resource also provides immediate infrastructure for reproducible questions, maintained answers, interface records, and event capture. Its broader contribution is a reproducible pattern for converting specialist technical discourse into transparent digital research objects whose claims remain bounded by their units, denominators, and evidence stages.
